% Created 2026-06-27 Sat 12:47
% Intended LaTeX compiler: pdflatex
\documentclass[11pt,letterpaper]{article}
\usepackage[utf8]{inputenc}
\usepackage[T1]{fontenc}
\usepackage{graphicx}
\usepackage{longtable}
\usepackage{wrapfig}
\usepackage{rotating}
\usepackage[normalem]{ulem}
\usepackage{amsmath}
\usepackage{amssymb}
\usepackage{capt-of}
\usepackage{hyperref}
\usepackage{amsmath}
\usepackage{amssymb}
\usepackage{geometry}
\geometry{left=1.4in, right=1.4in, top=1.2in, bottom=1.2in}
\usepackage{fancyhdr}
\usepackage{parskip}
\usepackage{microtype}
\usepackage{setspace}
\PassOptionsToPackage{hidelinks}{hyperref}
\setlength{\parindent}{0pt}
\setlength{\parskip}{8pt}
\onehalfspacing
\fancypagestyle{plain}{\fancyhf{}\fancyfoot[C]{\thepage}\renewcommand{\headrulewidth}{0pt}}
\fancypagestyle{body}{%
\fancyhf{}%
\fancyhead[C]{\small Document 0 of 5 $\cdot$ DOI: 10.5281/zenodo.18927154 $\cdot$ CC BY 4.0}%
\fancyfoot[C]{\thepage}%
\renewcommand{\headrulewidth}{0.4pt}}
\author{Bradley K. DePuy, Independent Researcher}
\date{2026}
\title{The Medium: A Personal Introduction}
\hypersetup{
 pdfauthor={Bradley K. DePuy, Independent Researcher},
 pdftitle={The Medium: A Personal Introduction},
 pdfkeywords={},
 pdfsubject={},
 pdfcreator={Emacs 29.3 (Org mode 9.6.15)}, 
 pdflang={English}}
\begin{document}

\pagestyle{plain}
\begin{titlepage}
\vspace*{4cm}
\begin{center}
  {\LARGE\bfseries The Medium}\\[0.6em]
  {\large An Attempted Description of a Pre-Geometric Substrate --- Personal Introduction}\\[2.5em]
  {\normalsize Bradley K. DePuy, Independent Researcher}\\[0.5em]
  {\normalsize ORCID: \mbox{\texttt{0009-0001-0328-2299}}}\\[0.5em]
  {\normalsize 2026}
\end{center}
\end{titlepage}
\pagestyle{body}

\section{Who I Am}
\label{sec:org0bfb8cb}

I have spent my life in software. Retired nearly 20 years ago and now have the time to explore questions plaguing me for some time.

I have wondered, what if I had pursued an academic path in Physics and Mathematics.  Probably both a blessing and a curse.
It would be a curse due to the constraint to prove my self within the tightly closed box of institutional acceptance.
This would have limited my ability to see other possibilities.
It would have been a blessing as I may now see the mathematical patterns necessary to expand the documents that follow.

\section{The Question}
\label{sec:orgcaf8299}

The question I have had since I don't remember, perhaps not specifically in these words, but expressed as:
--- Where does it all come from.

It is an existential question that may not have an answer.

Physics raises more questions than it solves.  It seems trapped in the institutional box mentioned above.  The beauty of the mathematics outweighs the truth behind what the math is really saying. Current physics must also agree with what is currently accepted and supposedly proven by observation.  I question whether those observations result from the evolution of what came before rather than an honest clean review of the assumptions.

The question itself has an assumption buried in it. "Why something rather than nothing" assumes that nothing is the natural state and something requires an explanation. But what if that is backwards? What if nothing is the harder thing to sustain? What if, given an infinite domain with no constraints and no preferred states, the emergence of something is not a puzzle requiring explanation but an inevitability? I cannot prove that. But it has always felt more honest than the alternatives --- more honest than a forced stop, more honest than a first cause that itself requires no cause.

I keep coming back to a handful of questions that the standard answers never quite settled for me. What if the cosmic background frequency is simply the stochastic noise of a Planck-scale substrate
rather than an echo of a singular event? And if that's possible, is the Hubble redshift really what we think it is, or have we borrowed a terrestrial concept and applied it somewhere it doesn't quite fit? Dark matter troubles me the same way — what if it isn't matter at all, but the energy signature of a substrate vibrating at scales we haven't learned to read yet? The mature galaxies turning up at the far boundaries feel like the same kind of clue. They shouldn't be there if everything started from a single point, but they make perfect sense if something has always been everywhere. Which brings me  to the Big Bang itself — maybe it happened, maybe it didn't, but perhaps it is one way of reading what a substrate that always was looks like when you observe it from inside.

What follows is my attempt to describe what might underlie it all, something prior to space and time, prior to particles.

The five documents that follow identify a number of unresolved and contested issues. I have not hidden them. An honest accounting of what works, what is uncertain, and what is simply wrong as stated is documented alongside the ideas themselves. I am continuing to dig into each of these problems and may resolve some of them. But I am realistic about what I am.

I am a sloppy philosopher. I can see patterns and ask questions that won't leave me alone. What I cannot do — at least not with the rigor these ideas deserve — is the disciplined mathematics needed to turn an attempted description into something testable and complete.

My hope is that someone finds this worth a second look. Perhaps a physicist who has wrestled with the same questions. Perhaps someone who hands it to a group of graduate students and says: here is an interesting set of ideas with real problems — see what you can do with it. The problems are documented. The structure is there. What it needs is someone with sharper tools than mine.

--- \emph{Bradley K. DePuy, 2026}
\end{document}
